\chapter{Tabella di comparazione tra NIS2, ISO/IEC 27001:2022 e NIST CSF 2.0}
\label{cap:mapping}

In questo capitolo si espone una tabella riassuntiva del confronto tra la Direttiva NIS2 e gli standard internazionali di cybersicurezza presi in esame.
Tale elaborato ha lo scopo di riassumere le considerazioni contenute nel Capitolo \ref{cap:confronto}, fornendo una visione più sintetica e operativa,
utile per la comprensione degli allineamenti e delle lacune tra la Direttiva, il framework NIST CSF 2.0 e lo standard ISO/IEC 27001:2022. 
In particolare, si mappano gli articoli della NIS2 con le corrispondenti misure previste negli obiettivi dell'allegato del framework NIST e nei controlli 
dell'\textit{Annex A} dello standard ISO. Si segnala che le corrispondenze qui individuate sono di carattere riassuntivo
e devono necessariamente essere relazionate alle considerazioni, più complete, riportate nel Capitolo \ref{cap:confronto}.

La tabella è costruita tramite un confronto diretto tra l'articolo 21 della Direttiva NIS2 e le relative misure specificate nel Regolamento di esecuzione 2024/2690,
lo standard ISO/IEC 27001:2022 e il framework NIST CSF 2.0. Sulla base dell’analisi condotta nel testo \cite{enisa-implementing-guidance}.

A supporto del confronto sono stati analizzati contributi similari presenti nelle seguenti fonti:
\begin{itemize}
    \item \cite{parmar-miles-csf};
    \item \cite{toussaint};
    \item \cite{mcintosh};
    \item \cite{uni-pdr174}.
\end{itemize}

\setlength{\tabcolsep}{4.5pt}          % padding orizzontale nelle celle
\renewcommand{\arraystretch}{1.4}      % padding verticale (fattore moltiplicativo)
\setlength{\extrarowheight}{1pt}       % padding verticale addizionale per riga

\setlength{\emergencystretch}{3em}     % riduce gli Underfull \hbox
\newcolumntype{P}[1]{>{\raggedright\arraybackslash}p{#1}} % colonne p raggedright

\section{Politica per la sicurezza delle reti e dei sistemi informativi}
\begin{longtable}{@{}|P{0.27\textwidth}|P{0.30\textwidth}|P{0.33\textwidth}|@{}} % Nota: assicurarsi che il preambolo carichi \usepackage{longtable}
\caption{Confronto tra NIS2, ISO/IEC 27001:2022 e NIST CSF 2.0 – Politica per la sicurezza delle reti e dei sistemi informativi}\label{tab:pol-sic-reti}\\
\hline
\rowcolor{gray!20}\textbf{NIS2} & \textbf{ISO/IEC 27001:2022} & \textbf{NIST CSF 2.0} \\ \hline
\endfirsthead
\hline
\rowcolor{gray!20}\textbf{NIS2} & \textbf{ISO/IEC 27001:2022} & \textbf{NIST CSF 2.0} \\ \hline
\endhead
21(2)(a): Politica per la sicurezza delle reti e dei sistemi informativi & 5.2, A.5.1, A.5.36, A.5.4, 9.3  & PR.AT-02, GV.PO-01, GV.PO-02, GV.OC-03, GV.RM-03, GV.OC-02, ID.IM-01, ID.IM-02, ID.IM-03, ID.IM-04 \\ \hline
21(2)(a): Ruoli, responsabilità e poteri correlati & 5.3, A.5.2, A.5.3, A.5.4  & GV.RR-02, GV.SC-02, PR.AT-02, ID.IM-01, ID.IM-02, ID.IM-03, ID.IM-04 \\ \hline
\end{longtable}

\section{Politica di gestione del rischio}
\begin{longtable}{@{}|P{0.27\textwidth}|P{0.30\textwidth}|P{0.33\textwidth}|@{}}
\caption{Confronto tra NIS2, ISO/IEC 27001:2022 e NIST CSF 2.0 – Politica di gestione del rischio}\label{tab:pol-gest-rischio}\\
\hline
\rowcolor{gray!20}\textbf{NIS2} & \textbf{ISO/IEC 27001:2022} & \textbf{NIST CSF 2.0} \\ \hline
\endfirsthead
\hline
\rowcolor{gray!20}\textbf{NIS2} & \textbf{ISO/IEC 27001:2022} & \textbf{NIST CSF 2.0} \\ \hline
\endhead
21(2)(a): Quadro per la gestione del rischio & 6.1, 6.1.2, 6.1.3, 6.2, 8.2, 8.3, A.5.7, A.5.19, A.5.20, A.5.21 & ID.RA-01, ID.RA-02, ID.RA-03, ID.RA-04, ID.RA-05, ID.RA-06, GV.RM-03, ID.RM-01, GV.RM-06, GV.RR-03, ID.IM-01, ID.IM-02, ID.IM-03, ID.IM-04 \\ \hline
21(2)(a): Monitoraggio della conformità & 9.2, A.5.31, A.5.35, A.5.36  & GV.OV-02, ID.IM-01, ID.IM-02, ID.IM-03, ID.IM-04 \\ \hline
21(2)(a): Riesame indipendente della sicurezza delle reti e delle informazioni & 9.2, 10.1, A.5.35, A.8.34  & GV.OV-02, ID.IM-01  \\ \hline
\end{longtable}

\section{Gestione degli incidenti}
\begin{longtable}{@{}|P{0.27\textwidth}|P{0.30\textwidth}|P{0.33\textwidth}|@{}}
\caption{Confronto tra NIS2, ISO/IEC 27001:2022 e NIST CSF 2.0 – Gestione degli incidenti}\label{tab:gest-incidenti}\\
\hline
\rowcolor{gray!20}\textbf{NIS2} & \textbf{ISO/IEC 27001:2022} & \textbf{NIST CSF 2.0} \\ \hline
\endfirsthead
\hline
\rowcolor{gray!20}\textbf{NIS2} & \textbf{ISO/IEC 27001:2022} & \textbf{NIST CSF 2.0} \\ \hline
\endhead
21(2)(b): Politica per la gestione degli incidenti & A.5.24  & GV.SC-08, RS.MA-01, RS.MA-05, RS.MI-01, RS.MI-02, ID.IM-01, ID.IM-04 \\ \hline
21(2)(b): Monitoraggio e registrazione (logging) & A.5.28, A.8.15, A.8.16, A.8.17 & RS.AN-06, RS.AN-07, ID.IM-01, ID.IM-02, ID.IM-03, ID.IM-04 \\ \hline
21(2)(b): Segnalazione degli eventi & A.6.8  & RS.MI-01, RS.CO-02  \\ \hline
21(2)(b): Valutazione e classificazione degli eventi & A.5.25  & DE.AE-04, RS.MA-02, RS.MA-03, RS.MA-04, ID.IM-01, ID.IM-02, ID.IM-03, ID.IM-04 \\ \hline
21(2)(b): Risposta agli incidenti & A.5.26  & RS.MA-01, RS.MA-02, RS.MA-03, RS.MA-04, ID.IM-02, ID.IM-03, ID.IM-04, RS.CO-02, RS.CO-03, RS.AN-03, RS.MI-01, RS.MI-02, RC.CO-03, RC.CO-04 \\ \hline
21(2)(b): Riesami post-incidente & A.5.27 & ID.IM-01, ID.IM-04, RS.AN-08 \\ \hline
\end{longtable}

\section{Continuità operativa e gestione delle crisi}
\begin{longtable}{@{}|P{0.27\textwidth}|P{0.30\textwidth}|P{0.33\textwidth}|@{}}
\caption{Confronto tra NIS2, ISO/IEC 27001:2022 e NIST CSF 2.0 – Continuità operativa e gestione delle crisi}\label{tab:continuita-crisi}\\
\hline
\rowcolor{gray!20}\textbf{NIS2} & \textbf{ISO/IEC 27001:2022} & \textbf{NIST CSF 2.0} \\ \hline
\endfirsthead
\hline
\rowcolor{gray!20}\textbf{NIS2} & \textbf{ISO/IEC 27001:2022} & \textbf{NIST CSF 2.0} \\ \hline
\endhead
21(2)(c): Piano di continuità operativa e ripristino in caso di disastro & A.5.29, A.5.30 & ID.IM-02, ID.IM-03, ID.IM-04, GV.OC-04, GV.SC-08, RC.RP-01, RC.RP-02 \\ \hline
21(2)(c): Gestione dei backup & A.8.13, A.8.14  & PR.DS-11, RC.RP-01, RC.RP-02, ID.IM-03 \\ \hline
21(2)(c): Gestione delle crisi & A.5.26, A.5.29, A.5.30  & RS.CO-02, RS.CO-03, PR.IR-03, DE.CM-01, ID.AM-03, DE.AE-02, DE.AE-03, DE.AE-04, DE.AE-06, DE.AE-07, DE.AE-08 \\ \hline
\end{longtable}

\section{Sicurezza della catena di fornitura}
\begin{longtable}{@{}|P{0.27\textwidth}|P{0.30\textwidth}|P{0.33\textwidth}|@{}}
\caption{Confronto tra NIS2, ISO/IEC 27001:2022 e NIST CSF 2.0 – Sicurezza della catena di fornitura}\label{tab:supply-chain}\\
\hline
\rowcolor{gray!20}\textbf{NIS2} & \textbf{ISO/IEC 27001:2022} & \textbf{NIST CSF 2.0} \\ \hline
\endfirsthead
\hline
\rowcolor{gray!20}\textbf{NIS2} & \textbf{ISO/IEC 27001:2022} & \textbf{NIST CSF 2.0} \\ \hline
\endhead
21(2)(d): Politica per la sicurezza della catena di fornitura & A.5.19, A.5.20, A.5.21, A.8.30 & GV.OC-03, GV.OC-05, GV.SC-01, GV.SC-04, GV.SC-06, GV.SC-05, GV.SC-07, GV.SC-09, GV.SC-10, ID.RA-10, ID.IM-01, ID.IM-02, ID.IM-03, ID.IM-04 \\ \hline
21(2)(d): Elenco di fornitori e prestatori di servizi & A.5.22  & GV.OC-05, GV.SC-04, ID.IM-01, ID.IM-02, ID.IM-03, ID.IM-04 \\ \hline
\end{longtable}

\section{Sicurezza nell’acquisizione, nello sviluppo e nella manutenzione di sistemi e servizi ICT}
\begin{longtable}{@{}|P{0.27\textwidth}|P{0.30\textwidth}|P{0.33\textwidth}|@{}}
\caption{Confronto tra NIS2, ISO/IEC 27001:2022 e NIST CSF 2.0 – Sicurezza nell’acquisizione, nello sviluppo e nella manutenzione di sistemi e servizi ICT}\label{tab:sviluppo-manunzione-ict}\\
\hline
\rowcolor{gray!20}\textbf{NIS2} & \textbf{ISO/IEC 27001:2022} & \textbf{NIST CSF 2.0} \\ \hline
\endfirsthead
\hline
\rowcolor{gray!20}\textbf{NIS2} & \textbf{ISO/IEC 27001:2022} & \textbf{NIST CSF 2.0} \\ \hline
\endhead
21(2)(e): Sicurezza di servizi, sistemi o prodotti ICT & A.5.21, A.5.23  & GV.PO-02, GV.SC-06, ID.RA-09, ID.RA-10, ID.IM-01, ID.IM-02, ID.IM-03, ID.IM-04 \\ \hline
21(2)(e): Ciclo di sviluppo sicuro (SDLC) & A.8.25, A.8.31 & ID.AM-08, PR.PS-06, ID.IM-01, ID.IM-02, ID.IM-03, ID.IM-04 \\ \hline
21(2)(e): Gestione della configurazione & A.8.9  & PR.PS-01, ID.IM-01, ID.IM-02, ID.IM-03, ID.IM-04 \\ \hline
21(2)(e): Gestione dei cambiamenti, riparazioni e manutenzione & 6.3, 8.1, A.7.13, A.8.32  & ID.RA-07, ID.IM-01, ID.IM-02, ID.IM-03, ID.IM-04 \\ \hline
21(2)(e): Verifiche di sicurezza & A.8.29, A.8.33, A.8.34  & ID.RA-01, ID.IM-01, ID.IM-02, ID.IM-03, ID.IM-04 \\ \hline
21(2)(e): Gestione delle patch di sicurezza & A.8.31, A.8.32  & PR.PS-02, DE.CM-09, ID.IM-01, ID.IM-02, ID.IM-03, ID.IM-04 \\ \hline
21(2)(e): Sicurezza di rete & A.8.16, A.8.20  & DE.CM-01, PR.IR-01, PR.PS-05, ID.IM-01, ID.IM-02, ID.IM-03, ID.IM-04 \\ \hline
21(2)(e): Segmentazione di rete & A.8.22  & PR.IR-01, ID.IM-01, ID.IM-02, ID.IM-03, ID.IM-04 \\ \hline
21(2)(e): Protezione da software dannoso e non autorizzato & A.5.32, A.8.7  & DE.CM-01, DE.CM-09, PR.PS-05, ID.IM-01, ID.IM-02, ID.IM-03, ID.IM-04 \\ \hline
21(2)(e): Gestione e divulgazione delle vulnerabilità & A.8.8 & ID.RA-01, ID.RA-02, ID.RA-04, ID.RA-05, ID.RA-06, PR.PS-02, PR.PS-03, ID.RA-08, ID.IM-01, ID.IM-02, ID.IM-03, ID.IM-04 \\ \hline
\end{longtable}

\section{Politiche e procedure per valutare l’efficacia delle misure di gestione del rischio di cybersicurezza}
\begin{longtable}{@{}|P{0.27\textwidth}|P{0.30\textwidth}|P{0.33\textwidth}|@{}}
\caption{Confronto tra NIS2, ISO/IEC 27001:2022 e NIST CSF 2.0 – Politiche e procedure per valutare l’efficacia delle misure di gestione del rischio di cybersicurezza}\label{tab:politiche-procedure-efficacia}\\
\hline
\rowcolor{gray!20}\textbf{NIS2} & \textbf{ISO/IEC 27001:2022} & \textbf{NIST CSF 2.0} \\ \hline
\endfirsthead
\hline
\rowcolor{gray!20}\textbf{NIS2} & \textbf{ISO/IEC 27001:2022} & \textbf{NIST CSF 2.0} \\ \hline
\endhead
21(2)(f): Politiche e procedure di valutazione dell’efficacia & 6.2, 9.1, 9.3  & ID.IM-03, GV.RM-06, ID.IM-01, ID.IM-02, ID.IM-03, ID.IM-04 \\ \hline
\end{longtable}

\section{Buone pratiche di igiene informatica e formazione sulla sicurezza}
\begin{longtable}{@{}|P{0.27\textwidth}|P{0.30\textwidth}|P{0.33\textwidth}|@{}}
\caption{Confronto tra NIS2, ISO/IEC 27001:2022 e NIST CSF 2.0 – Buone pratiche di igiene informatica e formazione sulla sicurezza}\label{tab:igiene-formazione}\\
\hline
\rowcolor{gray!20}\textbf{NIS2} & \textbf{ISO/IEC 27001:2022} & \textbf{NIST CSF 2.0} \\ \hline
\endfirsthead
\hline
\rowcolor{gray!20}\textbf{NIS2} & \textbf{ISO/IEC 27001:2022} & \textbf{NIST CSF 2.0} \\ \hline
\endhead
21(2)(g): Sensibilizzazione e buone pratiche di igiene informatica & 7.3, A.6.3, A.8.7  & PR.AT-01, PR.AT-02, ID.IM-01, ID.IM-02, ID.IM-03, ID.IM-04 \\ \hline
21(2)(g): Formazione sulla sicurezza & 7.2, A.6.3  & PR.AT-01, ID.IM-01, ID.IM-02, ID.IM-03, ID.IM-04 \\ \hline
\end{longtable}

\section{Crittografia}
\begin{longtable}{@{}|P{0.27\textwidth}|P{0.30\textwidth}|P{0.33\textwidth}|@{}}
\caption{Confronto tra NIS2, ISO/IEC 27001:2022 e NIST CSF 2.0 – Crittografia}\label{tab:crittografia}\\
\hline
\rowcolor{gray!20}\textbf{NIS2} & \textbf{ISO/IEC 27001:2022} & \textbf{NIST CSF 2.0} \\ \hline
\endfirsthead
\hline
\rowcolor{gray!20}\textbf{NIS2} & \textbf{ISO/IEC 27001:2022} & \textbf{NIST CSF 2.0} \\ \hline
\endhead
21(2)(h): Crittografia & A.5.31, A.8.24  & PR.DS-01, PR.DS-02, ID.IM-01, ID.IM-02, ID.IM-03, ID.IM-04 \\ \hline
\end{longtable}

\section{Sicurezza delle risorse umane}
\begin{longtable}{@{}|P{0.27\textwidth}|P{0.30\textwidth}|P{0.33\textwidth}|@{}}
\caption{Confronto tra NIS2, ISO/IEC 27001:2022 e NIST CSF 2.0 – Sicurezza delle risorse umane}\label{tab:risorse-umane}\\
\hline
\rowcolor{gray!20}\textbf{NIS2} & \textbf{ISO/IEC 27001:2022} & \textbf{NIST CSF 2.0} \\ \hline
\endfirsthead
\hline
\rowcolor{gray!20}\textbf{NIS2} & \textbf{ISO/IEC 27001:2022} & \textbf{NIST CSF 2.0} \\ \hline
\endhead
21(2)(i): Sicurezza delle risorse umane & 7.1, 7.2, A.6.2, A.6.3  & PR.AT-02, GV.RR-04, ID.IM-01, ID.IM-02, ID.IM-03, ID.IM-04 \\ \hline
21(2)(i): Verifica dei precedenti & A.6.1  & GV.RR-04, ID.IM-01, ID.IM-02, ID.IM-03, ID.IM-04 \\ \hline
21(2)(i): Procedure di cessazione o variazione del rapporto di lavoro & A.6.5  & GV.RR-04  \\ \hline
21(2)(i): Procedimento disciplinare & A.6.4  & ID.IM-01, ID.IM-02, ID.IM-03, ID.IM-04 \\ \hline
\end{longtable}

\section{Controllo degli accessi}
\begin{longtable}{@{}|P{0.27\textwidth}|P{0.30\textwidth}|P{0.33\textwidth}|@{}}
\caption{Confronto tra NIS2, ISO/IEC 27001:2022 e NIST CSF 2.0 – Controllo degli accessi}\label{tab:controllo-accessi}\\
\hline
\rowcolor{gray!20}\textbf{NIS2} & \textbf{ISO/IEC 27001:2022} & \textbf{NIST CSF 2.0} \\ \hline
\endfirsthead
\hline
\rowcolor{gray!20}\textbf{NIS2} & \textbf{ISO/IEC 27001:2022} & \textbf{NIST CSF 2.0} \\ \hline
\endhead
21(2)(i): Politica di controllo degli accessi & A.5.15, A.7.2, A.8.3, A.8.21, A.9 & PR.AA-05, ID.IM-01, ID.IM-02, ID.IM-03, ID.IM-04 \\ \hline
21(2)(i): Gestione dei diritti di accesso & A.5.18, A.9  & PR.AA-05, ID.IM-01, ID.IM-02, ID.IM-03, ID.IM-04 \\ \hline
21(2)(i): Account privilegiati e account di amministrazione dei sistemi & A.8.2, A.8.18, A.9  & ID.IM-01, ID.IM-02, ID.IM-03, ID.IM-04 \\ \hline
21(2)(i): Sistemi di amministrazione & A.8.2, A.8.18  & Nessuna corrispondenza diretta \\ \hline
21(2)(i): Identificazione & A.5.16  & PR.AA-01, PR.AA-05, PR.AC-02 \\ \hline
21(2)(i): Autenticazione & A.5.17  & PR.AA-05, PR.AA-03, ID.IM-01, ID.IM-02, ID.IM-03, ID.IM-04 \\ \hline
21(2)(i): Autenticazione a più fattori & A.8.5  & PR.AA-03  \\ \hline
\end{longtable}

\section{Gestione degli asset}
\begin{longtable}{@{}|P{0.27\textwidth}|P{0.30\textwidth}|P{0.33\textwidth}|@{}}
\caption{Confronto tra NIS2, ISO/IEC 27001:2022 e NIST CSF 2.0 – Gestione degli asset}\label{tab:gestione-asset}\\
\hline
\rowcolor{gray!20}\textbf{NIS2} & \textbf{ISO/IEC 27001:2022} & \textbf{NIST CSF 2.0} \\ \hline
\endfirsthead
\hline
\rowcolor{gray!20}\textbf{NIS2} & \textbf{ISO/IEC 27001:2022} & \textbf{NIST CSF 2.0} \\ \hline
\endhead
21(2)(i): Classificazione degli asset & A.5.9, A.5.12, A.5.13  & ID.AM-05  \\ \hline
21(2)(i): Trattamento degli asset & A.5.9, A.5.10, A.5.14, A.7.10  & ID.IM-01  \\ \hline
21(2)(i): Politica sui supporti rimovibili & A.7.7, A.7.10  & PR.DS-01, ID.IM-01, ID.IM-02, ID.IM-03, ID.IM-04 \\ \hline
21(2)(i): Inventario degli asset & A.5.9  & ID.AM-01, ID.AM-02, ID.AM-03, ID.AM-04, ID.AM-07, ID.AM-08, ID.IM-01, ID.IM-02, ID.IM-03, ID.IM-04 \\ \hline
21(2)(i): Consegna, restituzione o cancellazione degli asset alla cessazione del rapporto di lavoro & A.5.11, A.5.18, A.8.24  & Nessuna corrispondenza diretta \\ \hline
\end{longtable}

\section{Sicurezza fisica e ambientale}
\begin{longtable}{@{}|P{0.27\textwidth}|P{0.30\textwidth}|P{0.33\textwidth}|@{}}
\caption{Confronto tra NIS2, ISO/IEC 27001:2022 e NIST CSF 2.0 – Sicurezza fisica e ambientale}\label{tab:sic-fisica-ambientale}\\
\hline
\rowcolor{gray!20}\textbf{NIS2} & \textbf{ISO/IEC 27001:2022} & \textbf{NIST CSF 2.0} \\ \hline
\endfirsthead
\hline
\rowcolor{gray!20}\textbf{NIS2} & \textbf{ISO/IEC 27001:2022} & \textbf{NIST CSF 2.0} \\ \hline
\endhead
21(2)(c): Servizi di supporto (utilities) & A.7.11  & DE.CM-02, DE.CM-06, GV.OC-03, GV.OC-05, GV.OC-07, ID.RA-10, ID.IM-01, ID.IM-02, ID.IM-03, ID.IM-04 \\ \hline
21(2)(e): Protezione dalle minacce fisiche e ambientali & A.7.3, A.7.5  & PR.IR-02, ID.IM-01, ID.IM-02, ID.IM-03, ID.IM-04 \\ \hline
21(2)(i): Perimetro e controllo degli accessi fisici & A.7.1, A.7.2, A.7.4  & PR.AA-06, DE.CM-02  \\ \hline
\end{longtable}